\section{Resultados y comparación}

La respuesta natural deberá contrastarse con la frecuencia amortiguada y con el factor de decaimiento $e^{-\zeta\omega_nt}$. Para la respuesta forzada se verificará que el desplazamiento parta de cero, responda mientras actúa la perturbación y regrese al equilibrio una vez desaparezca la fuerza, siempre que los polos se encuentren en el semiplano izquierdo.

La comparación considerará, como mínimo, la amplitud máxima, el número de oscilaciones visibles, el tiempo de decaimiento y el comportamiento final. También deberá distinguirse la respuesta provocada por energía inicial de aquella producida por la fuerza externa.

\draftnote{Insertar aquí las figuras exportadas por Julia, sus parámetros, las verificaciones numéricas y la interpretación física. Cada figura debe incluir título, ejes, unidades y leyenda.}

